\titre{}
\theme{calcLit}
\auteur{Nathan Scheinmann}
\niveau{1M}
\source{nathan}
\type{serie}
\piments{4}
\pts{}
\annee{2425}

\contenu{
\tcblower
	Exprimer la variable demandée en fonction des autres variables présentes dans la formule.
\begin{tasks}
	\task $V=\dfrac{h}{6}\left(B_1+B_2+4M\right) \hfill h=? \hfill M=? \hfill$
	\task $D_r=\dfrac{D}{1+A_rB_r} \hfill D=? \hfill A_r=? \hfill$
	\task $P=Q\dfrac{R-r}{2R} \hfill r=? \hfill \phantom{1} \hfill$
	\task $G=\dfrac{kR_a}{R_i+R_a} \hfill R_i=? \hfill \phantom{1} \hfill$
	\task $A=\dfrac{F+S_\alpha}{S_\alpha} \hfill F=? \hfill \phantom{1} \hfill$
	\task $\dfrac{1}{R}=\dfrac{1}{R_1}+\dfrac{1}{R_2} \hfill R=? \hfill R_1=? \hfill$
\end{tasks}	
}
\correction{
\tcblower
\begin{enumerate}
	\item $V=\dfrac{h}{6}\left(B_1+B_2+4M\right) \hfill h=? \hfill M=? \hfill$
		
		Isolons $h$:

		\begin{solve}
			V&=\dfrac{h}{6}(B_1+B_2+4M)&\div (B_1+B_2+4M)\\
			\dfrac{V}{B_1+B_2+4M)}&=\dfrac{h}{6}&\cdot 6\\
			\dfrac{6V}{B_1+B_2+4M)}&=h&h\text{ est isolé}
		\end{solve}
		Isolons $M$:
		\begin{solve}
			V&=\dfrac{h}{6}(B_1+B_2+4M)&\cdot \dfrac{6}{h}\\
			\dfrac{6V}{h}&=B_1+B_2+4M&-B_1-B_2\\
			\dfrac{6V}{h}-B_1-B_2&=4M&\div 4\\
			\dfrac{\dfrac{6V}{h}-B_1-B_2}{4}&=M&M \text{ est isolé}
		\end{solve}

	\item $D_r=\dfrac{D}{1+A_rB_r} \hfill D=? \hfill A_r=? \hfill$

		Isolons $D$:
		\begin{solve}
			D_r&=\dfrac{D}{1+A_r+B_r}&\cdot (1+A_r+B_r)\\
			D_r\cdot (1+A_r+B_r)&=D&D \text{ est isolé}
		\end{solve}
	
		Isolons $A_r$ (on reprend la dernière formule du point précédent):
		\begin{solve}
			D_r\cdot (1+A_r+B_r)&=D&\div D_r\\
			1+A_r+B_r&=\dfrac{D}{D_r}&-1-B_r\\
			A_r&=\dfrac{D}{D_r}-1-B_r&A_r\text{ est isolé}
		\end{solve}
	\item $P=Q\dfrac{R-r}{2R} \hfill r=? \hfill \phantom{1} \hfill$
	
		Isolons $r$:
		\begin{solve}
			P&=Q\dfrac{R-r}{2R}&\div Q\\
			\dfrac{P}{Q}&=\dfrac{R-r}{2R}&\cdot 2R\\
			\dfrac{P\cdot 2R}{Q}&=R-r&-R\\
			\dfrac{2PR}{Q}&=-r&\cdot (-1)\\
			-\dfrac{2PR}{Q}&=r&r \text{ est isolé}
		\end{solve}

	\item $G=\dfrac{kR_a}{R_i+R_a} \hfill R_i=? \hfill \phantom{1} \hfill$
		
		Isolons $R_i$:
		\begin{solve}
			G&=\dfrac{kR_a}{R_i+R_a}&\cdot (R_i+R_a)\\
			G(R_i+R_a)&=kR_a&\div G\\
			R_i+R_a&=\dfrac{kR_a}{G}&-R_a\\
			R_i&=\dfrac{kR_a}{G}-R_a&R_a\text{ est isolé}
		\end{solve}

	\item $A=\dfrac{F+S_\alpha}{S_\alpha} \hfill F=? \hfill \phantom{1} \hfill$

		Isolons $F$:
		\begin{solve}
			A&=\dfrac{F+S_\alpha}{S_{\alpha}}&\cdot S_\alpha\\
			\dfrac{A}{S_\alpha}&=F+S_\alpha&-S_\alpha\\
			\dfrac{A}{S_\alpha}-S_\alpha&=F&F \text{ est isolé}
		\end{solve}
	\item $\dfrac{1}{R}=\dfrac{1}{R_1}+\dfrac{1}{R_2} \hfill R=? \hfill R_1=? \hfill$
		
		On commence par écrire le membre de droite sous une autre forme:
		\begin{align*}
			\dfrac{1}{R_1}+\dfrac{1}{R_2}&=\dfrac{R_2}{R_1R_2}+\dfrac{R_1}{R_1R_2}
						     &=\dfrac{R_2+R_1}{R_1R_2}
		\end{align*}

		Isolons $R$:
		\begin{solve}
		\dfrac{1}{R}&=\dfrac{1}{R_1}+\dfrac{1}{R_2}&\text{Autre forme}\\
			\dfrac{1}{R}&=\dfrac{R_2+R_1}{R_1R_2}&\cdot R\\
			1&=\dfrac{R_2+R_1}{R_1R_2}\cdot R&\cdot \dfrac{R_1R_2}{R_1+R_2}\\
			\dfrac{R_1R_2}{R_1+R_2}&=R& R \text{ est isolé}
		\end{solve}

		Isolons $R_1$:
		\begin{solve}
			\dfrac{1}{R}&=\dfrac{1}{R_1}+\dfrac{1}{R_2}&-\dfrac{1}{R_2}\\
			\dfrac{1}{R}-\dfrac{1}{R_2}&=\dfrac{1}{R_1}&\\
			\dfrac{R_2-R}{RR_2}&=\dfrac{1}{R_1}&\cdot R_1\\
			\dfrac{R_2-R}{RR_2}\cdot R_1&=1&\cdot \dfrac{RR_2}{R_2-R}\\
			R_1&=\dfrac{RR_2}{R_2-R}&R_1\text{ est isolé}

		\end{solve}
\end{enumerate}	

}

